\documentclass[11pt]{article}
\usepackage[T1]{fontenc}
\usepackage[margin=1.1in]{geometry}
\usepackage{booktabs}
\usepackage{longtable}
\usepackage{array}
\usepackage{listings}
\usepackage{xcolor}
\usepackage[table]{xcolor}
\usepackage{geometry}
\geometry{margin=1in}
\definecolor{utilcolor}{rgb}{0.9,1.0,0.9}
\usepackage{hyperref}
\hypersetup{colorlinks=true, linkcolor=blue!60!black}
\usepackage{microtype}
\usepackage{parskip}
\usepackage{enumitem}
\emergencystretch=3em

% ── JSON listing style ────────────────────────────────────────────────────────
\definecolor{jsonstring}{RGB}{0,120,0}
\definecolor{jsonnumber}{RGB}{180,90,0}
\definecolor{jsonbool}{RGB}{0,140,140}
\definecolor{codebg}{RGB}{250,250,250}
\definecolor{promptbg}{RGB}{245,245,255}
\lstdefinelanguage{json}{
    basicstyle=\ttfamily\footnotesize,
    backgroundcolor=\color{codebg},
    showstringspaces=false,
    breaklines=true,
    frame=single, framesep=4pt,
    rulecolor=\color{gray!40},
    morestring=[b]",
    stringstyle=\color{jsonstring},
    literate=
        *{0}{{{\color{jsonnumber}0}}}{1}
         {1}{{{\color{jsonnumber}1}}}{1}
         {2}{{{\color{jsonnumber}2}}}{1}
         {3}{{{\color{jsonnumber}3}}}{1}
         {4}{{{\color{jsonnumber}4}}}{1}
         {5}{{{\color{jsonnumber}5}}}{1}
         {6}{{{\color{jsonnumber}6}}}{1}
         {7}{{{\color{jsonnumber}7}}}{1}
         {8}{{{\color{jsonnumber}8}}}{1}
         {9}{{{\color{jsonnumber}9}}}{1}
         {true}{{{\color{jsonbool}true}}}{4}
         {false}{{{\color{jsonbool}false}}}{5}
         {null}{{{\color{gray}null}}}{4},
}
\lstdefinestyle{prompt}{
    basicstyle=\ttfamily\footnotesize,
    backgroundcolor=\color{promptbg},
    showstringspaces=false,
    breaklines=true,
    frame=single, framesep=4pt,
    rulecolor=\color{gray!40},
}

\newcommand{\app}[1]{\texttt{#1}}

\begin{document}

% =============================================================================
\section{Synthetic Application Suite}
\label{sec:synthetic_apps}
% =============================================================================

To broaden the application coverage of AppWorld beyond its original eleven productivity-focused applications, we construct a suite of 52 synthetic applications spanning ten functional categories. Each application is modeled as a REST API consistent with AppWorld's structural conventions and is fully specified with typed endpoints, parameter constraints, and structured response schemas. The suite comprises 894 API endpoints across the 50 core applications (mean 17.9 per application, range 15--21).

% --
\subsection{Domain Selection}
% --

The 52 application domains were selected through a manual curation process guided by two principles: \emph{real-world fidelity}, where each domain corresponds to a widely used category of software, and \emph{functional diversity}, where the collection as a whole requires agents to exercise qualitatively different interaction patterns, data models, and reasoning skills. Domains were selected to complement and extend AppWorld's existing application coverage while minimizing redundancy with previously supported domains.

We manually defined a seed set of 10 broad functional categories and identified specific app types within each category. For each app type we specified a short description and a list of key agent-facing functionalities. Listing~\ref{lst:seed_domain} shows a representative entry from our seed set. Using these manually defined entries as in-context examples, we then prompted Gemini to expand the list with additional domains following the same format and coverage criteria, specifically requesting diversity across the remaining domain space and avoiding overlap with the seed entries.

\begin{lstlisting}[language=json,
  caption={A representative entry from the manually curated seed domain set used to prompt Gemini for further domain expansion. Each entry specifies the \texttt{general\_domain}, the specific \texttt{name}, a \texttt{description}, and a comma-separated list of agent \texttt{functionalities}.},
  label={lst:seed_domain}]
{
  "general_domain": "Finance & Payments",
  "name": "Peer-to-Peer Money Transfer",
  "description": "A digital wallet app that lets users send and receive
      money, split bills, and manage payment methods.",
  "functionalities": "send money, request payments, view transaction
      history, manage payment cards, add/withdraw balance,
      like and comment on transactions"
}
\end{lstlisting}

The resulting 52 applications are listed in Table~\ref{tab:apps}, grouped by functional category. Two applications, \texttt{inboxly} (email client, 48 APIs) and \texttt{shopnow} (e-commerce, 75 APIs), were added subsequently to fill identified coverage gaps and have larger API surfaces than the core 50 apps; they are not included in the suite total of 894.

Table~\ref{tab:domains} summarises the ten functional domains and their constituent app counts.

% --
\subsubsection{Email and Shopping Coverage: \texttt{inboxly} and \texttt{shopnow}}
% --

After generating the core 50 apps, we identified two significant coverage gaps: the suite contained no email client and no general e-commerce platform. These are two of the richest domains in AppWorld Test-C (\texttt{gmail} and \texttt{amazon}), and leaving them absent would limit the diversity of multi-app tasks. We therefore generated \texttt{inboxly} and \texttt{shopnow} as independent synthetic analogues using \texttt{gemini-2.5-pro}.

For each app, we provided five AppWorld API endpoints as structural few-shot examples, selected to cover five distinct functional areas of the target domain. For \texttt{inboxly} we sampled from the AppWorld \texttt{gmail} schema; for \texttt{shopnow} we sampled from the AppWorld \texttt{amazon} schema. These five examples anchor Gemini to the target domain's vocabulary and interaction patterns while the prompt explicitly instructs it to generate a new, independent app with different naming and structure.

\begin{lstlisting}[language=json,
  caption={Five AppWorld \texttt{gmail} APIs provided as structural examples when generating \texttt{inboxly}, covering five distinct functional areas: authentication, thread listing with rich filters, email sending with attachments, thread labelling, and cross-app file download.},
  label={lst:inboxly_examples}]
1. login              POST /gmail/auth/token
   params: username, password
   returns: { "access_token": "string", "token_type": "string" }

2. show_inbox_threads GET /gmail/email_threads/category/inbox
   params: access_token, query, label, starred, archived, spam,
           snoozed, read, attachment, from_email, to_email,
           min_created_at, max_created_at, page_index, page_limit, sort_by
   returns: [{ "email_thread_id": 1, "subject": "string", ... }]

3. send_email         POST /gmail/emails
   params: email_addresses (req), subject (req), body (req),
           access_token, attachment_file_paths, file_system_access_token
   returns: { "message": "string", "sent_email_thread_id": 1 }

4. label_thread       POST /gmail/email_threads/{email_thread_id}/label
   params: email_thread_id (req), label (req), access_token
   returns: { "message": "string" }

5. download_attachment GET /gmail/attachments/{attachment_id}
   params: attachment_id (req), file_system_access_token (req),
           access_token, download_to_file_path, overwrite
   returns: { "message": "string", "file_path": "string" }
\end{lstlisting}

\begin{lstlisting}[language=json,
  caption={Five AppWorld \texttt{amazon} APIs provided as structural examples when generating \texttt{shopnow}, covering five distinct functional areas: authentication, product search with rich filters, cart management, order placement, and review submission.},
  label={lst:shopnow_examples}]
1. login              POST /amazon/auth/token
   params: username, password
   returns: { "access_token": "string", "token_type": "string" }

2. search_products    GET /amazon/products
   params: query, product_type, color, relative_size,
           min_price, max_price, min_product_rating, max_product_rating,
           min_seller_rating, max_seller_rating, seller_id,
           page_index, page_limit, sort_by
   returns: [{ "product_id": 1, "name": "string", "price": 0.0, ... }]

3. add_product_to_cart POST /amazon/cart/{product_id}
   params: product_id (req), access_token,
           quantity (>= 1), clear_cart_first
   returns: { "message": "string" }

4. place_order        POST /amazon/orders
   params: payment_card_id (req), address_id (req), access_token
   returns: { "message": "string", "order_id": 1 }

5. write_product_review POST /amazon/products/{product_id}/reviews
   params: product_id (req), rating (req, 1-5), access_token,
           title, text
   returns: { "message": "string", "product_review_id": 1 }
\end{lstlisting}

The prompt then specified the target domain explicitly. For \texttt{inboxly}, the target header read:

\begin{quote}\ttfamily\small
Your job is to invent ONE new \textbf{Personal Email Client} app (in the \textbf{Social \& Communication} category). App context: A personal email client for managing inbox, sending emails, and organizing with labels. Key workflows to cover: send and receive emails, manage inbox/sent/starred/archived/spam threads, handle drafts, manage labels, search messages, forward and reply, manage attachments. IMPORTANT: Invent a new, distinct app name — do not reuse any well-known brand name.
\end{quote}

For \texttt{shopnow}, the target header read:

\begin{quote}\ttfamily\small
Your job is to invent ONE new \textbf{General E-Commerce} app (in the \textbf{E-Commerce \& Marketplace} category). App context: A general-purpose online shopping marketplace. Key workflows to cover: browse and search products, manage cart and wishlist, place orders, track deliveries, write reviews, manage returns, handle payment methods and addresses. IMPORTANT: Invent a new, distinct app name — do not reuse any well-known brand name.
\end{quote}

\paragraph{inboxly.}
\texttt{inboxly} is a personal email client with 48 APIs covering the full email workflow: account management and authentication, inbox and folder navigation (inbox, sent, starred, archived, spam, snoozed threads), thread and message operations (read, delete, label, mark), draft management (create, edit, send, attach files), and contact search. Its API surface is deliberately independent from \texttt{gmail}: rather than \texttt{show\_inbox\_threads}, \texttt{inboxly} exposes \texttt{inbox\_thread\_list}; rather than \texttt{mark\_thread\_starred}, it uses \texttt{highlight\_thread}; and its draft lifecycle uses \texttt{composition} terminology throughout.

\paragraph{shopnow.}
\texttt{shopnow} is a general e-commerce platform with 75 APIs covering: account management, product browsing and search, seller discovery, browsing history, cart and wishlist management (add, update, move between cart and wishlist), gift wrapping, order placement and receipts, payment card and address management, product reviews and Q\&A, return management, and membership/subscription plans. Its naming conventions differ systematically from \texttt{amazon}: \texttt{add\_item\_to\_cart} instead of \texttt{add\_product\_to\_cart}, \texttt{checkout} instead of \texttt{place\_order}, \texttt{get\_item\_reviews} instead of \texttt{show\_product\_reviews}, and \texttt{membership\_tiers} instead of Prime plans.

\begin{table}[h]
\centering
\small
\caption{The ten functional domains of the synthetic app suite, obtained by manual inspection of real-world software categories. Each domain was defined to represent a distinct interaction paradigm for LLM tool-use agents.}
\label{tab:domains}
\setlength{\tabcolsep}{6pt}
\begin{tabular}{lcp{8.5cm}}
\toprule
\textbf{Domain} & \textbf{\#Apps} & \textbf{Representative interaction patterns} \\
\midrule
Social \& Communication      & 7 & Post, follow, message, react, moderate content \\
Productivity \& Proj.\ Mgmt  & 7 & Create tasks, manage boards, schedule events, sign docs \\
Entertainment \& Media       & 7 & Stream, subscribe, rate, bookmark, manage playlists \\
E-Commerce \& Marketplace    & 5 & Browse, bid, purchase, review, manage listings \\
Travel \& Transportation     & 5 & Book, track, navigate, estimate fares, manage trips \\
Food \& Health               & 5 & Order food, log meals, track fitness, manage health records \\
Developer \& Business Tools  & 5 & Manage repos, file tickets, run campaigns, track leads \\
Finance \& Payments          & 4 & Transfer funds, trade assets, manage accounts, pay bills \\
Local \& Community           & 3 & Discover venues, RSVP to events, join groups \\
Learning \& Hobbies          & 3 & Complete lessons, track workouts, check weather \\
\midrule
\textbf{Total}               & \textbf{52} & \\
\bottomrule
\end{tabular}
\end{table}

\setlength{\LTcapwidth}{\linewidth}
{\small
\setlength{\tabcolsep}{4pt}
\renewcommand{\arraystretch}{0.85}
\begin{longtable}{@{}>{\bfseries\raggedright\arraybackslash}p{2.65cm} >{\raggedright\arraybackslash}p{2.55cm} >{\raggedright\arraybackslash}p{8.30cm} r@{}}
\caption{The 52 synthetic applications grouped by functional domain, sorted by descending API count within each domain. \texttt{inboxly} and \texttt{shopnow} were added after the main pipeline run and have larger API surfaces than the core 50 apps.}
\label{tab:apps} \\
\toprule
\textbf{Domain} & \textbf{App} & \textbf{Description} & \textbf{\#} \\
\midrule
\endfirsthead
\multicolumn{4}{l}{\small\textit{Table~\ref{tab:apps} continued}} \\[1pt]
\toprule
\textbf{Domain} & \textbf{App} & \textbf{Description} & \textbf{\#} \\
\midrule
\endhead
\midrule
\multicolumn{4}{r}{\small\textit{Continued on next page}} \\
\endfoot
\bottomrule
\endlastfoot

Social \& Comm.
  & \app{chirper}    & Microblogging network for short messages, following users, and trending topics.          & 21 \\
  & \app{teamhive}   & Team messaging app for channels, direct messages, file sharing, and reminders.            & 20 \\
  & \app{snapgrid}   & Photo and video social network for sharing posts, commenting, and direct messaging.        & 19 \\
  & \app{proconnect} & Professional networking for connecting with colleagues, posting updates, and job search.  & 19 \\
  & \app{inkpost}    & Long-form writing platform for publishing articles and managing a reading list.            & 19 \\
  & \app{forumhub}   & Community discussion forum for posting, voting, and subscribing to subforums.              & 18 \\
  & \app{inboxly}    & Personal email client for managing inbox, sending emails, and organizing with labels.      & 48 \\
\midrule

Productivity \& Proj.\ Mgmt
  & \app{boardflow}  & Project management tool with boards, lists, and cards for team collaboration.             & 20 \\
  & \app{bugtrack}   & Issue tracking system for managing bugs, sprints, epics, and backlogs.                    & 19 \\
  & \app{pageflow}   & Workspace and note-taking app for pages, databases, and collaborative documents.           & 19 \\
  & \app{clouddrop}  & Cloud file storage for uploading, organizing, sharing, and managing files.                 & 19 \\
  & \app{taskboard}  & Task management for organizing projects, assigning tasks, and team collaboration.          & 18 \\
  & \app{meetspace}  & Video conferencing for scheduling meetings, polls, recordings, and webinars.               & 18 \\
  & \app{daysync}    & Calendar and scheduling app for events, attendees, and shared calendars.                   & 16 \\
  & \app{signseal}   & Digital document signing for creating envelopes and tracking signature status.             & 16 \\
\midrule

Entertainment \& Media
  & \app{viewtube}    & Video sharing platform for uploading, watching, liking, and subscribing to channels.      & 20 \\
  & \app{streamvault} & Video streaming service for browsing, watching, and downloading movies and TV shows.      & 18 \\
  & \app{audiowave}   & Music streaming for uploading, discovering, liking, and sharing tracks and playlists.     & 18 \\
  & \app{bookshelf}   & Social book-tracking for cataloging books, writing reviews, and rating reads.             & 18 \\
  & \app{newsflash}   & Personalized news aggregator for following topics, reading articles, and saving content.  & 17 \\
  & \app{podlisten}   & Podcast app for subscribing to shows, playing episodes, and managing playlists.           & 16 \\
  & \app{pageturn}    & Digital bookstore and e-reader for purchasing, reading, and annotating ebooks.            & 16 \\
\midrule

Travel \& Transportation
  & \app{swiftride}   & Ride-hailing app for requesting rides, tracking drivers, and managing payments.           & 18 \\
  & \app{journeybook} & Travel booking for flights, hotels, and car rentals, and managing trips.                 & 17 \\
  & \app{homeseek}    & Real estate search for finding properties, contacting agents, and mortgage estimates.     & 17 \\
  & \app{pathfinder}  & Navigation app for directions, nearby places, traffic, and saved locations.               & 17 \\
  & \app{travelpal}   & Travel discovery and planning for destinations, hotels, and restaurants.                  & 16 \\
\midrule

E-Commerce \& Marketplace
  & \app{nestshare}   & Short-term rental marketplace for booking accommodations and managing stays.              & 19 \\
  & \app{craftmarket} & Handmade goods marketplace for buying and selling crafts and managing orders.            & 18 \\
  & \app{bidmart}     & Online auction marketplace for placing bids, buying items, and creating listings.         & 17 \\
  & \app{locallist}   & Classifieds marketplace for jobs, housing, items for sale, and services.                  & 16 \\
  & \app{shopnow}     & E-commerce platform for browsing products, managing cart, and placing orders.             & 75 \\
\midrule

Food \& Health
  & \app{quickbite}  & Food delivery app for browsing restaurants, placing orders, and tracking deliveries.      & 21 \\
  & \app{medrecord}  & Personal health records for conditions, medications, allergies, and appointments.         & 20 \\
  & \app{nutrilog}   & Nutrition and fitness tracker for logging meals, exercises, and health goals.             & 18 \\
  & \app{pulsetrack} & Fitness wearable companion for activity, sleep, heart rate, and health goals.             & 18 \\
  & \app{tablebook}  & Restaurant reservation app for finding restaurants and booking tables.                    & 16 \\
\midrule

Developer \& Business
  & \app{codevault}    & Code hosting for managing repositories, branches, and pull requests.                    & 20 \\
  & \app{leadflow}     & CRM for managing contacts, deals, emails, tasks, and sales pipeline analytics.          & 19 \\
  & \app{helpdesk}     & Customer support ticketing for creating, assigning, and resolving tickets.              & 18 \\
  & \app{mailblast}    & Email marketing for managing subscriber lists, campaigns, and bulk sends.               & 17 \\
  & \app{messagecloud} & Cloud communications API for sending SMS messages, calls, and phone numbers.           & 16 \\
\midrule

Finance \& Payments
  & \app{vaultbank}  & Online banking for checking balances, transferring funds, and paying bills.               & 19 \\
  & \app{cryptonest} & Cryptocurrency portfolio tracker for buying, selling, and monitoring digital assets.      & 19 \\
  & \app{paylink}    & Peer-to-peer payments for sending money, invoices, and managing subscriptions.            & 18 \\
  & \app{stocknest}  & Online brokerage for buying and selling stocks and tracking a portfolio.                  & 18 \\
\midrule

Local \& Community
  & \app{localspot} & Local business discovery for searching restaurants, reading reviews, and saving favorites. & 16 \\
  & \app{gatherup}  & Local groups and events app for communities, RSVPs, and group messaging.                  & 16 \\
  & \app{eventpass} & Event ticketing for creating events, selling tickets, and managing registrations.         & 16 \\
\midrule

Learning \& Hobbies
  & \app{trailblazer} & Fitness and outdoor activity tracker for workouts, segments, and athlete leaderboards.  & 18 \\
  & \app{lingualeap}  & Language learning app for courses, lessons, streaks, and leaderboards.                   & 17 \\
  & \app{skywatch}    & Weather forecast app for conditions, hourly forecasts, alerts, and air quality.           & 15 \\

\end{longtable}
}

% --
\subsection{API Schema Design}
% --

For each domain, we generated a complete REST API schema by prompting Gemini (\texttt{gemini-2.5-pro}) with (1) the target domain description drawn from our domain taxonomy, and (2) six representative Venmo endpoints from AppWorld as structural few-shot examples. The six examples were deliberately selected to cover the full range of structural patterns that appear in the suite: a public authentication endpoint (\texttt{login}), a simple authenticated \texttt{GET} (\texttt{show\_account}), a resource creation \texttt{POST} with required and optional parameters (\texttt{create\_transaction}), a complex list endpoint with filters, pagination, and \texttt{sort\_by} (\texttt{show\_transactions}), a stateful action on a sub-resource (\texttt{like\_transaction}), and a sub-resource \texttt{GET} by ID (\texttt{show\_payment\_card}). Listing~\ref{lst:example_apis} shows two of these six examples.

\begin{lstlisting}[language=json,
  caption={Two of the six Venmo APIs provided as few-shot structural examples to Gemini: a complex list endpoint with filters and pagination (top) and a stateful sub-resource action (bottom). These two examples illustrate the range of structural patterns Gemini is expected to replicate.},
  label={lst:example_apis}]
{
  "app_name": "venmo",  "api_name": "show_transactions",
  "path": "/venmo/transactions",  "method": "GET",
  "description": "Search or show a list of your transactions.",
  "parameters": [
    {"name": "access_token", "type": "string",  "required": true,
     "default": null, "constraints": []},
    {"name": "query",        "type": "string",  "required": false,
     "default": "", "constraints": []},
    {"name": "min_amount",   "type": "number",  "required": false,
     "default": 0.01, "constraints": ["value > 0.0"]},
    {"name": "direction",    "type": "string",  "required": false,
     "default": null,
     "constraints": ["value in ['sent', 'received']"]},
    {"name": "page_index",   "type": "integer", "required": false,
     "default": 0,    "constraints": ["value >= 0"]},
    {"name": "page_limit",   "type": "integer", "required": false,
     "default": 5,    "constraints": ["value >= 1, <= 20"]},
    {"name": "sort_by",      "type": "string",  "required": false,
     "default": null, "constraints": []}
  ],
  "response_schemas": {
    "success": [{"transaction_id": 1, "amount": 0.0,
      "description": "string", "created_at": "2019-01-01T00:00:00",
      "private": true,
      "sender":   {"name": "string", "email": "user@example.com"},
      "receiver": {"name": "string", "email": "user@example.com"}}],
    "failure": {"message": "string"}
  }
},
{
  "app_name": "venmo",  "api_name": "like_transaction",
  "path": "/venmo/transactions/{transaction_id}/likes",
  "method": "POST",
  "description": "Like a transaction.",
  "parameters": [
    {"name": "transaction_id", "type": "integer", "required": true,
     "default": null, "constraints": []},
    {"name": "access_token",   "type": "string",  "required": true,
     "default": null, "constraints": []}
  ],
  "response_schemas": {
    "success": {"message": "string"},
    "failure": {"message": "string"}
  }
}
\end{lstlisting}

\paragraph{Prompt structure.}
The full Gemini prompt is structured in five parts, shown in abbreviated form in Listing~\ref{lst:gemini_prompt}. The domain name and description from the domain taxonomy (Listing~\ref{lst:seed_domain}) are injected at the top of the prompt as a \emph{target header} (Part~1), which tells Gemini which specific domain to generate and which key workflows to cover. Part~(2) specifies the required \texttt{app\_description} and the API coverage requirements (core CRUD, sub-resources, actions, search/filter, social features, notifications, payments, file downloads). Part~(3) defines the exact JSON output schema for each API object and each parameter object. Part~(4) states mandatory conventions: the \texttt{access\_token} parameter ordering rule, the 10 standard auth/account APIs every app must include, pagination conventions, and response schema placeholder formats. Part~(5) provides the six Venmo example APIs. A diversity constraint is appended when generating multiple apps in sequence, listing already-generated app names that must not be duplicated.

\begin{lstlisting}[style=prompt,
  caption={Abbreviated structure of the Gemini prompt used for synthetic app generation. Domain entries are loaded from \texttt{domains\_subdomains.json} at runtime; \texttt{<domain.*>} placeholders are filled from the current entry. Ellipses (\texttt{...}) denote sections paraphrased for brevity.},
  label={lst:gemini_prompt}]
# Load domain entries from file
domains = json.load(open("domains_subdomains.json"))

# For each domain entry, build and send this prompt:
# --- Part 1: Target header (domain entry injected from domains_subdomains.json) ---
Your job is to invent ONE new <domain.name> app
(in the <domain.general_domain> category).

App context: <domain.description>
Key workflows to cover: <domain.functionalities>

IMPORTANT: Invent a new, distinct app name.
           Do not reuse any well-known brand name.

# --- Part 2: Coverage requirements ---
You are an expert API designer building a new mobile app backend
from scratch. The app should support at least 30-50 APIs.

Beyond the 10 standard account management APIs, cover:
  1. Core entities: CRUD for each major resource
  2. Sub-resources: nested resources (e.g. order items, comments)
  3. Actions: stateful transitions (e.g. cancel_order, approve_request)
  4. Search & filtering: list APIs with filters, date ranges, enums
  5. Social features: follows, likes, ratings, reviews, comments
  6. Notifications: show_notifications, mark_notification, delete_notification
  7. Payments: payment methods, charges, refunds, receipts
  8. File downloads: receipts/invoices following venmo download_*_receipt

# --- Part 3: Output format ---
A JSON object { "<app_name>": { "app_description": "...", "apis": [...] }}
Each API: { "app_name", "api_name", "path", "method",
  "description", "parameters", "response_schemas" }
Each parameter: { "name", "type", "required", "description",
  "default", "constraints" }

# --- Part 4: Conventions ---
access_token: required on all non-public APIs; order is
  required params -> access_token -> optional params.
Public APIs (no access_token): login, signup,
  send_verification_code, send_password_reset_code,
  reset_password, verify_account.
10 required standard APIs: signup, login, logout, show_account,
  update_account, send_verification_code, verify_account,
  send_password_reset_code, reset_password, delete_account.
Pagination on all list APIs: page_index (default 0),
  page_limit (default 5), sort_by (optional).
Response placeholders: "string", 1, 0.0, true,
  "2019-01-01T00:00:00", "user@example.com".
Constraints: "length >= 1", "value >= 0",
  "value in ['a','b']", "value is email address".

# --- Part 5: Few-shot examples ---
## Example: 6 representative Venmo APIs
[... six full API objects as shown in Listing 2 ...]

# --- Diversity constraint (appended when N > 1) ---
The following apps have already been generated. Your new app MUST
be from a completely different domain: `<app1>`, `<app2>`, ...

Output only the JSON object, no explanation, no markdown.
\end{lstlisting}

\paragraph{Verification.}
After each Gemini response, a rule-based verifier checks 13 structural invariants before accepting the schema: all required top-level keys are present; all 10 standard APIs are included; each endpoint has the seven required fields; paths start with \texttt{/<app\_name>/}; HTTP methods are valid; \texttt{access\_token} is absent on public APIs and present on all others; \texttt{access\_token} does not appear after optional parameters; \texttt{response\_schemas} has both \texttt{success} and \texttt{failure} keys; all parameter types are from the allowed set; and at least one domain-specific API exists beyond the standard 10. Schemas that fail verification are discarded and regenerated up to a configurable retry limit.

\paragraph{Schema structure.}
Each endpoint is defined by eight fields: \texttt{app\_name}, \texttt{api\_name}, \texttt{path}, \texttt{method} (\texttt{GET}/\texttt{POST}/\texttt{PATCH}/\texttt{DELETE}), \texttt{description}, \texttt{parameters}, \texttt{response\_schemas}, and \texttt{canary\_string}. Specifications are authored using four Python helpers that enforce structural uniformity across all applications: \texttt{api\_()} constructs a full endpoint entry; \texttt{p()} constructs a typed parameter; \texttt{tok()} is shorthand for the standard \texttt{access\_token} parameter; and \texttt{paged()} appends the three standard pagination parameters to a parameter list.

\paragraph{Design conventions.}
We enforced the following conventions across all 52 applications to ensure a consistent interface vocabulary:

\begin{itemize}[leftmargin=1.5em, itemsep=2pt]
    \item \textbf{Authentication.} Every application exposes a \texttt{login} endpoint returning an \texttt{access\_token}, which is required on all authenticated endpoints, mirroring the pattern in the original AppWorld applications.
    \item \textbf{Pagination.} All list-returning endpoints accept \texttt{page\_index} (0-based, default 0), \texttt{page\_limit} (default 5), and \texttt{sort\_by} (attribute prefixed with \texttt{+}/\texttt{-} for ascending/descending order).
    \item \textbf{Type-safe response schemas.} Response schemas use typed placeholders: \texttt{"string"} for text fields, \texttt{1} for integer IDs, \texttt{0.0} for floats, \texttt{"2019-01-01T00:00:00"} for timestamps, and \texttt{"user@example.com"} for email fields, never domain-specific invented values.
    \item \textbf{Constraint notation.} Parameter constraints follow a strict vocabulary: \texttt{"length >= 1"}, \texttt{"value >= 0.0, <= 5.0"}, \texttt{"value in ['a', 'b']"}, and \texttt{"value is email address"}. No other formats are valid.
    \item \textbf{Nested sub-resources.} Where a resource logically contains an embedded sub-resource (e.g., a ride has a driver; a booking has a host), the response schema includes a nested object rather than flattening all fields.
    \item \textbf{Failure responses.} The \texttt{failure} branch of every endpoint is uniformly \texttt{\{"message": "string"\}} across all 52 applications.
\end{itemize}

% \paragraph{Example: \texttt{swiftride}.}
% The \texttt{swiftride} application models a ride-hailing service. Its 18 endpoints cover the complete ride lifecycle: authentication, account management, fare estimation, ride dispatch, real-time tracking, cancellation, rating, driver discovery, payment methods, promotion codes, and issue reporting. Listing~\ref{lst:swiftride} shows three representative endpoints: an authentication endpoint, a single-resource \texttt{GET} with a nested sub-object, and a paginated list with an enum filter, illustrating the three dominant structural patterns present throughout the suite.

\begin{lstlisting}[language=json,
  caption={Three endpoints from the \texttt{swiftride} schema, illustrating authentication (top), a single-resource \texttt{GET} with a nested \texttt{driver} object (middle), and a paginated list with an enum filter (bottom). Canary strings are truncated.},
  label={lst:swiftride}]
{
  "login": {
    "app_name": "swiftride",  "api_name": "login",
    "path": "/swiftride/login",  "method": "POST",
    "description": "Log in to swiftride and obtain an access token.",
    "parameters": [
      {"name": "email",    "type": "string", "required": true,
       "default": null, "constraints": ["value is email address"]},
      {"name": "password", "type": "string", "required": true,
       "default": null, "constraints": ["length >= 1"]}
    ],
    "response_schemas": {
      "success": {"access_token": "string", "message": "string"},
      "failure": {"message": "string"}
    },
    "canary_string": "appworld:3ee6fbe:8d128d36-..."
  },
  "show_ride": {
    "app_name": "swiftride",  "api_name": "show_ride",
    "path": "/swiftride/rides/{ride_id}",  "method": "GET",
    "description": "Show details of a specific ride.",
    "parameters": [
      {"name": "ride_id",      "type": "integer", "required": true},
      {"name": "access_token", "type": "string",  "required": true}
    ],
    "response_schemas": {
      "success": {
        "ride_id": 1, "pickup_address": "string",
        "dropoff_address": "string", "fare": 0.0,
        "status": "string", "created_at": "2019-01-01T00:00:00",
        "driver": {"name": "string", "rating": 0.0,
                   "vehicle": "string", "license_plate": "string"}
      },
      "failure": {"message": "string"}
    },
    "canary_string": "appworld:3ee6fbe:c8f13ef5-..."
  },
  "show_rides": {
    "app_name": "swiftride",  "api_name": "show_rides",
    "path": "/swiftride/rides",  "method": "GET",
    "description": "Show your ride history with optional status filter.",
    "parameters": [
      {"name": "access_token", "type": "string",  "required": true},
      {"name": "status", "type": "string", "required": false,
       "default": null,
       "constraints": ["value in ['completed','cancelled','scheduled']"]},
      {"name": "page_index", "type": "integer", "default": 0},
      {"name": "page_limit", "type": "integer", "default": 5},
      {"name": "sort_by",    "type": "string",  "default": null}
    ],
    "response_schemas": {
      "success": [{"ride_id": 1, "pickup_address": "string",
        "dropoff_address": "string", "fare": 0.0,
        "status": "string", "created_at": "2019-01-01T00:00:00"}],
      "failure": {"message": "string"}
    },
    "canary_string": "appworld:3ee6fbe:fa291c3a-..."
  }
}
\end{lstlisting}

% --
\subsection{Schema Conversion}
% --

The Python specification files are serialized to a consolidated JSON file and then converted per application into AppWorld's canonical format using Gemini (\texttt{gemini-3.1-pro-preview}). For each application, a structured prompt is constructed containing: (1) reference examples from the original AppWorld Spotify schema as few-shot demonstrations; (2) explicit formatting rules for description style, constraint vocabulary, and response schema placeholder conventions; and (3) the full API specification for the target application. Gemini returns a JSON object mapping each \texttt{api\_name} to its complete AppWorld-format entry.

Following conversion, each API entry is annotated with a \emph{canary string} of the form \texttt{appworld:\{md5\_prefix\}:\{uuid\}}, where the MD5 prefix is the first seven hex characters of the application name's hash and the UUID is freshly generated per endpoint. Canary strings serve two purposes: they enable traceability of downstream generated data back to its source application, and they provide a signal for detecting schema memorization during model evaluation. The conversion pipeline includes automatic retry on JSON parse failures and falls back to a deterministic structural converter for any applications where LLM conversion fails, ensuring all 50 applications produce valid output.

% --
\subsection{AppWorld Integration}
% --

The converted schemas are consumed by a deterministic code generation script that produces a deployable AppWorld module for each application, comprising five files: \texttt{\_\_init\_\_.py} (module initialization), \texttt{info.toml} (metadata with \texttt{app\_type = "synthetic"} marker), \texttt{models.py} (SQLModel database models with a standard \texttt{User} table), \texttt{apis.py} (FastAPI endpoint stubs wired to AppWorld's shared authentication library), and \texttt{responses.py} (Pydantic response model definitions). The 50 synthetic applications are registered as \texttt{SYNTHETIC\_APP\_NAMES} separately from the 11 original \texttt{ORIGINAL\_APP\_NAMES}, and the \texttt{app\_type = "synthetic"} marker enables controlled ablation studies over application subsets during evaluation.




\section{API Comparison: Our Synthetic Apps vs.\ AppWorld Test-C Apps}

\noindent Our synthetic apps operate in the same functional domains as AppWorld Test-C apps
(shopping, email) but are independently implemented with distinct API interfaces.
While some domain-specific vocabulary necessarily overlaps (any shopping app will reference
concepts like carts, payments, and reviews; any email app will reference threads and drafts),
the apps differ in API signatures, parameter structures, response schemas, and the multi-step
workflow patterns required to complete tasks. An agent trained only on AppWorld Test-N apps
must therefore learn to navigate a new API surface at test time, rather than recalling
memorized call patterns from training.

\subsection*{Shopping: \texttt{shopnow} (ours) vs.\ \texttt{amazon} (AppWorld Test-C)}

\colorbox{utilcolor}{\strut Utility APIs (both apps)}\\[6pt]

\begin{longtable}{ll}
\caption{Comparison of shopping APIs between our synthetic app \texttt{shopnow} and AppWorld's
\texttt{amazon}. Green rows denote utility APIs present in both apps;
remaining rows list each app's functional APIs independently.}\\
\toprule
\textbf{shopnow API} & \textbf{amazon API} \\
\endfirsthead
\caption[]{(continued)}\\
\toprule
\textbf{shopnow API} & \textbf{amazon API} \\
\endhead
\bottomrule
\endfoot
\midrule
\multicolumn{2}{c}{\cellcolor{utilcolor}\small\textit{Utility APIs (common to both apps)}} \\
\midrule
\multicolumn{2}{c}{\cellcolor{utilcolor}\texttt{sign\_in}} \\
\multicolumn{2}{c}{\cellcolor{utilcolor}\texttt{sign\_out}} \\
\multicolumn{2}{c}{\cellcolor{utilcolor}\texttt{register}} \\
\multicolumn{2}{c}{\cellcolor{utilcolor}\texttt{view\_account}} \\
\multicolumn{2}{c}{\cellcolor{utilcolor}\texttt{view\_profile}} \\
\multicolumn{2}{c}{\cellcolor{utilcolor}\texttt{edit\_account\_name}} \\
\multicolumn{2}{c}{\cellcolor{utilcolor}\texttt{change\_password}} \\
\multicolumn{2}{c}{\cellcolor{utilcolor}\texttt{request\_password\_reset}} \\
\multicolumn{2}{c}{\cellcolor{utilcolor}\texttt{request\_verification}} \\
\multicolumn{2}{c}{\cellcolor{utilcolor}\texttt{confirm\_account}} \\
\multicolumn{2}{c}{\cellcolor{utilcolor}\texttt{close\_account}} \\
\midrule
\multicolumn{2}{l}{\small\textit{App-specific APIs}} \\
\midrule
\texttt{get\_payment\_options} & \texttt{show\_product} \\
\texttt{open\_return\_case} & \texttt{show\_recommended\_products} \\
\texttt{add\_payment\_option} & \texttt{show\_browsing\_history} \\
\texttt{clear\_shopping\_cart} & \texttt{clear\_browsing\_history} \\
\texttt{get\_question\_responses} & \texttt{add\_product\_to\_browsing\_history} \\
\texttt{get\_item\_details} & \texttt{remove\_product\_from\_browsing\_history} \\
\texttt{get\_last\_purchase} & \texttt{update\_browsing\_history\_tracking} \\
\texttt{get\_order\_details} & \texttt{show\_last\_product\_purchase} \\
\texttt{update\_item\_question} & \texttt{show\_product\_rating\_distribution} \\
\texttt{browse\_catalog} & \texttt{search\_sellers} \\
\texttt{update\_cart\_item\_quantity} & \texttt{show\_seller} \\
\texttt{download\_purchase\_receipt} & \texttt{search\_product\_types} \\
\texttt{update\_wishlist\_item\_quantity} & \texttt{show\_product\_feature\_choices} \\
\texttt{submit\_item\_review} & \texttt{search\_products} \\
\texttt{delete\_question\_response} & \texttt{show\_cart} \\
\texttt{delete\_item\_review} & \texttt{clear\_cart} \\
\texttt{update\_item\_review} & \texttt{add\_product\_to\_cart} \\
\texttt{get\_rating\_breakdown} & \texttt{update\_product\_quantity\_in\_cart} \\
\texttt{move\_wishlist\_item\_to\_cart} & \texttt{delete\_product\_from\_cart} \\
\texttt{search\_merchants} & \texttt{apply\_promo\_code\_to\_cart} \\
\texttt{update\_question\_response} & \texttt{remove\_promo\_code\_from\_cart} \\
\texttt{add\_item\_to\_wishlist} & \texttt{show\_wish\_list} \\
\texttt{download\_membership\_receipt} & \texttt{clear\_wish\_list} \\
\texttt{get\_purchased\_items} & \texttt{add\_product\_to\_wish\_list} \\
\texttt{add\_item\_to\_history} & \texttt{delete\_product\_from\_wish\_list} \\
\texttt{get\_payment\_option\_details} & \texttt{update\_product\_quantity\_in\_wish\_list} \\
\texttt{delete\_item\_question} & \texttt{move\_product\_from\_cart\_to\_wish\_list} \\
\texttt{get\_item\_filters} & \texttt{move\_product\_from\_wish\_list\_to\_cart} \\
\texttt{update\_shipping\_address} & \texttt{add\_gift\_wrapping\_to\_product} \\
\texttt{update\_history\_tracking} & \texttt{remove\_gift\_wrapping\_from\_product} \\
\texttt{basket\_overview} & \texttt{show\_orders} \\
\texttt{update\_payment\_option} & \texttt{place\_order} \\
\texttt{join\_membership\_plan} & \texttt{show\_product\_purchases} \\
\texttt{get\_item\_questions} & \texttt{show\_order} \\
\texttt{get\_wishlist\_items} & \texttt{download\_order\_receipt} \\
\texttt{remove\_item\_from\_wishlist} & \texttt{show\_payment\_cards} \\
\texttt{add\_item\_to\_cart} & \texttt{add\_payment\_card} \\
\texttt{get\_membership\_plans} & \texttt{show\_payment\_card} \\
\texttt{delete\_shipping\_address} & \texttt{update\_payment\_card} \\
\texttt{get\_item\_reviews} & \texttt{delete\_payment\_card} \\
\texttt{remove\_gift\_wrap\_from\_item} & \texttt{show\_addresses} \\
\texttt{clear\_item\_history} & \texttt{add\_address} \\
\texttt{get\_merchant\_details} & \texttt{update\_address} \\
\texttt{add\_shipping\_address} & \texttt{delete\_address} \\
\texttt{remove\_item\_from\_cart} & \texttt{show\_product\_reviews} \\
\texttt{keyword\_search} & \texttt{write\_product\_review} \\
\texttt{submit\_question\_response} & \texttt{update\_product\_review} \\
\texttt{recent\_views} & \texttt{delete\_product\_review} \\
\texttt{personalized\_picks} & \texttt{show\_product\_questions} \\
\texttt{delivery\_locations} & \texttt{write\_product\_question} \\
 & \texttt{show\_product\_question\_answers} \\
 & \texttt{write\_product\_question\_answer} \\
 & \texttt{update\_product\_question} \\
 & \texttt{delete\_product\_question} \\
 & \texttt{update\_product\_question\_answer} \\
 & \texttt{delete\_product\_question\_answer} \\
 & \texttt{show\_returns} \\
 & \texttt{initiate\_return} \\
 & \texttt{show\_return} \\
 & \texttt{show\_return\_deliverers} \\
 & \texttt{show\_prime\_plans} \\
 & \texttt{subscribe\_prime} \\
 & \texttt{show\_prime\_subscriptions} \\
 & \texttt{download\_prime\_subscription\_receipt} \\
\end{longtable}

\subsection*{Email: \texttt{inboxly} (ours) vs.\ \texttt{gmail} (AppWorld Test-C)}

\colorbox{utilcolor}{\strut Utility APIs (both apps)}\\[6pt]

\begin{longtable}{ll}
\caption{Comparison of email APIs between our synthetic app \texttt{inboxly} and AppWorld's
\texttt{gmail}. Green rows denote utility APIs present in both apps;
remaining rows list each app's functional APIs independently.}\\
\toprule
\textbf{inboxly API} & \textbf{gmail API} \\
\endfirsthead
\caption[]{(continued)}\\
\toprule
\textbf{inboxly API} & \textbf{gmail API} \\
\endhead
\bottomrule
\endfoot
\midrule
\multicolumn{2}{c}{\cellcolor{utilcolor}\small\textit{Utility APIs (common to both apps)}} \\
\midrule
\multicolumn{2}{c}{\cellcolor{utilcolor}\texttt{sign\_in}} \\
\multicolumn{2}{c}{\cellcolor{utilcolor}\texttt{sign\_out}} \\
\multicolumn{2}{c}{\cellcolor{utilcolor}\texttt{register}} \\
\multicolumn{2}{c}{\cellcolor{utilcolor}\texttt{view\_account}} \\
\multicolumn{2}{c}{\cellcolor{utilcolor}\texttt{view\_profile}} \\
\multicolumn{2}{c}{\cellcolor{utilcolor}\texttt{edit\_account\_name}} \\
\multicolumn{2}{c}{\cellcolor{utilcolor}\texttt{change\_password}} \\
\multicolumn{2}{c}{\cellcolor{utilcolor}\texttt{request\_password\_reset}} \\
\multicolumn{2}{c}{\cellcolor{utilcolor}\texttt{find\_users}} \\
\multicolumn{2}{c}{\cellcolor{utilcolor}\texttt{update\_status}} \\
\multicolumn{2}{c}{\cellcolor{utilcolor}\texttt{close\_account}} \\
\midrule
\multicolumn{2}{l}{\small\textit{App-specific APIs}} \\
\midrule
\texttt{get\_draft\_item} & \texttt{search\_labels} \\
\texttt{get\_mail\_thread} & \texttt{show\_inbox\_threads} \\
\texttt{remove\_message\_entry} & \texttt{show\_outbox\_threads} \\
\texttt{unflag\_junk\_thread} & \texttt{show\_snoozed\_threads} \\
\texttt{highlighted\_thread\_list} & \texttt{show\_starred\_threads} \\
\texttt{send\_draft\_item} & \texttt{show\_archived\_threads} \\
\texttt{exclude\_file} & \texttt{show\_spam\_threads} \\
\texttt{highlight\_thread} & \texttt{show\_category\_sizes} \\
\texttt{composition\_list} & \texttt{show\_thread} \\
\texttt{discard\_draft\_item} & \texttt{delete\_thread} \\
\texttt{remove\_label\_from\_thread} & \texttt{show\_email} \\
\texttt{revise\_composition} & \texttt{label\_thread} \\
\texttt{label\_list} & \texttt{unlabel\_thread} \\
\texttt{restore\_conversation} & \texttt{mark\_thread\_read} \\
\texttt{get\_folder\_counts} & \texttt{mark\_thread\_unread} \\
\texttt{undefer\_mail\_thread} & \texttt{mark\_thread\_archived} \\
\texttt{respond\_in\_thread} & \texttt{mark\_thread\_unarchived} \\
\texttt{get\_message\_content} & \texttt{mark\_thread\_spam} \\
\texttt{download\_attached\_file} & \texttt{mark\_thread\_not\_spam} \\
\texttt{sent\_thread\_list} & \texttt{mark\_thread\_starred} \\
\texttt{flag\_as\_unread} & \texttt{mark\_thread\_unstarred} \\
\texttt{junk\_thread\_list} & \texttt{delete\_email\_in\_thread} \\
\texttt{unhighlight\_thread} & \texttt{snooze\_thread} \\
\texttt{mail\_feed} & \texttt{unsnooze\_thread} \\
\texttt{new\_draft\_item} & \texttt{send\_email} \\
\texttt{store\_conversation} & \texttt{reply\_to\_email} \\
\texttt{flag\_as\_read} & \texttt{forward\_email\_from\_thread} \\
\texttt{include\_files} & \texttt{forward\_email\_thread} \\
 & \texttt{create\_draft} \\
 & \texttt{show\_drafts} \\
 & \texttt{update\_draft} \\
 & \texttt{delete\_draft} \\
 & \texttt{show\_draft} \\
 & \texttt{send\_email\_from\_draft} \\
 & \texttt{download\_attachment} \\
 & \texttt{upload\_attachments\_to\_draft} \\
 & \texttt{remove\_attachment\_from\_draft} \\
\end{longtable}

\end{document}
